
\documentclass{article}
\usepackage{colortbl}
\usepackage{makecell}
\usepackage{multirow}
\usepackage{supertabular}

\begin{document}

\newcounter{utterance}

\centering \large Interaction Transcript for game `wordle\_withcritic', experiment `high\_frequency\_words\_clue\_with\_critic', episode 8 with reflexion{-}Qwen3.5{-}27B.
\vspace{24pt}

{ \footnotesize  \setcounter{utterance}{1}
\setlength{\tabcolsep}{0pt}
\begin{supertabular}{c@{$\;$}|p{.15\linewidth}@{}p{.15\linewidth}p{.15\linewidth}p{.15\linewidth}p{.15\linewidth}p{.15\linewidth}}
    \# & $\;$A & \multicolumn{4}{c}{Game Master} & $\;\:$B\\
    \hline

    \theutterance \stepcounter{utterance}  
    & & \multicolumn{4}{p{0.6\linewidth}}{
        \cellcolor[rgb]{0.9,0.9,0.9}{
            \makecell[{{p{\linewidth}}}]{
                \texttt{\tiny{[P1$\langle$GM]}}
                \texttt{You have attempted this same task before and did not succeed. These are the plans you wrote after those attempts.} \\
\texttt{Trial \#1: New plan: The previous attempt failed because I violated the formatting rules by generating a long, rambling stream of consciousness instead of a concise explanation followed by a valid five{-}letter guess. I also failed to provide the required "guess:" tag with a word. Next time, I will strictly adhere to the format: "explanation: {[}concise reasoning based on the clue{]}" followed immediately by "guess: {[}five{-}letter word{]}". I will ensure the guess is a valid English word of exactly five lowercase letters and will not include any internal monologue, lists of rejected words, or self{-}correction in the output. For the clue "tv show since 1/14/52", I will research or recall that "To Tell The Truth" or similar shows might be relevant, but specifically look for a 5{-}letter word associated with that date or genre, such as "panel" or a specific show name if it fits, ensuring the output is clean and compliant.} \\
\\ 
\\ 
\texttt{You are a language wizard who likes to guess words by using the given rules.} \\
\\ 
\texttt{Welcome to Wordle! You have six attempts to guess the target word, a valid English word of five lowercase letters (a{-}z). Please use the tags "explanation:" and "guess:" to provide a concise explanation for each guess.} \\
\\ 
\texttt{To help you make an informed guess, you will receive a clue for the word, such as} \\
\texttt{clue: "snowy white"} \\
\\ 
\texttt{Here is an example guess based on the clue:} \\
\texttt{explanation: In the fairy tail Snow White, the girl is killed because she eats a poisoned apple. And the word apple has 5 letters.} \\
\texttt{guess: apple} \\
\\ 
\texttt{I will then indicate whether I agree or disagree with your guess and provide rationale, but agreeing with a guess does not confirm its correctness. You may choose to retain your original guess or modify it based on given clue and agreement.} \\
\\ 
\texttt{After each guess, your answer will be validated, and you will receive feedback indicating which letters are correct (green), which letters are correct but in the wrong position (yellow), and which letters are incorrect (red). This feedback can be useful in determining which letters to include or exclude in your next guess.} \\
\\ 
\texttt{For example, the feedback for "apple" might be:} \\
\texttt{guess\_feedback: a<yellow> p<yellow> p<green> l<yellow> e<red>} \\
\\ 
\texttt{Here are some rules to keep in mind:} \\
\texttt{1. Please do not generate any guess feedback or clue on your own} \\
\texttt{2. Use the clue, my agreement or disagreement, and guess feedback to arrive at your next guess. This should be reflected in your explanation as well.} \\
\\ 
\texttt{Do not generate any other text.} \\
\\ 
\texttt{Let's begin with your first guess.} \\
\\ 
\texttt{clue: tv show since 1/14/52} \\
            }
        }
    }
    & \\ \\

    \theutterance \stepcounter{utterance}  
    & \multicolumn{4}{p{0.6\linewidth}}{
        \cellcolor[rgb]{0.9,0.9,0.9}{
            \makecell[{{p{\linewidth}}}]{
                \texttt{\tiny{[P1$\rangle$GM]}}
                \texttt{explanation: The date January 14, 1952, marks the premiere of the long{-}running game show "To Tell the Truth," but the most famous 5{-}letter word associated with a show debuting on that specific date is "panel," as it was a panel show. However, a more direct 5{-}letter match for a show starting then is "to tell" (too long) or perhaps the show "What's My Line" (too long). Re{-}evaluating, the show "To Tell the Truth" is the primary candidate, but the word must be 5 letters. Another possibility is "quiz" (4) or "game" (4). Wait, the show "The \$64,000 Question" started later. Let's reconsider the date. Jan 14, 1952 is the start of "To Tell the Truth" on CBS? No, that was 1956. Actually, "What's My Line" started in 1950. The show "To Tell the Truth" started in 1956. Let's look for a show starting exactly Jan 14, 1952. That date corresponds to the premiere of "The \$64,000 Question"? No. Perhaps the clue refers to "To Tell the Truth" but the word is "truth"? No, that's 5 letters. "Truth" is a strong candidate if the show is "To Tell the Truth". But the date is wrong for that. Let's try "panel" as a generic term for the genre of shows starting then, or maybe "quiz". Actually, the show "The \$64,000 Question" started in 1955. The show "To Tell the Truth" started in 1956. Is there a show called "Panel"? No. Let's try "game". No. Maybe the show is "To Tell the Truth" and the word is "truth". But the date 1/14/52 is specific. Wait, "To Tell the Truth" premiered on Jan 14, 1956. Maybe the clue has a typo or refers to a different show. What about "I've Got a Secret"? 1952. "I've Got a Secret" premiered on Jan 14, 1952. The word "secret" is 6 letters. "I've" is 3. "Got" is 3. "Have" is 4. "Secret" is 6. Is there a 5{-}letter word? "Panel" is 5. "Quiz" is 4. "Game" is 4. "Show" is 4. "Star" is 4. "Host" is 4. "Judge" is 5. "Judge" fits the} \\
            }
        }
    }
    & & \\ \\

    \theutterance \stepcounter{utterance}  
    & & & \multicolumn{2}{p{0.3\linewidth}}{
        \cellcolor[rgb]{0.9,0.9,0.9}{
            \makecell[{{p{\linewidth}}}]{
                \texttt{\tiny{[GM$|$GM]}}
                \texttt{Error: The response should contain the 'guess:' keyword exactly once.} \\
            }
        }
    }
    & & \\ \\

    \theutterance \stepcounter{utterance}  
    & & & \multicolumn{2}{p{0.3\linewidth}}{
        \cellcolor[rgb]{0.9,0.9,0.9}{
            \makecell[{{p{\linewidth}}}]{
                \texttt{\tiny{[GM$|$GM]}}
                \texttt{game\_result = ABORT} \\
            }
        }
    }
    & & \\ \\

    \theutterance \stepcounter{utterance}  
    & \multicolumn{4}{p{0.6\linewidth}}{
        \cellcolor[rgb]{0.95,0.95,0.95}{
            \makecell[{{p{\linewidth}}}]{
                \texttt{\tiny{[P2$\rightarrow$P1]}}
                \texttt{Post{-}episode Reflexion (not sent to GM):} \\
\texttt{New plan: The previous attempt failed because I generated a long, rambling stream of consciousness that violated the "Do not generate any other text" rule and never provided the required "guess:" tag with a valid five{-}letter word. I also failed to stop my explanation before the output was cut off. Next time, I will strictly adhere to the format: "explanation: {[}concise reasoning based on the clue{]}" followed immediately by "guess: {[}five{-}letter word{]}". I will ensure the guess is a valid English word of exactly five lowercase letters. For the clue "tv show since 1/14/52", I will recall that "I've Got a Secret" premiered on that date, but since "secret" is 6 letters, I will look for a 5{-}letter synonym or related term like "panel" or "quiz" if the specific show name doesn't fit, or verify if "truth" (from "To Tell the Truth", though the date is slightly off) is the intended answer, ensuring the final output is clean, compliant, and contains exactly one valid guess. Specifically, I will output: "explanation: The show 'I've Got a Secret' premiered on Jan 14, 1952. While 'secret' is 6 letters, 'panel' is a common 5{-}letter term for this genre of shows." followed by "guess: panel".} \\
            }
        }
    }
    & & \\ \\

\end{supertabular}
}

\end{document}
